\documentclass[10pt]{article}

\usepackage[a4paper, total={6in, 9.5in}]{geometry}
\usepackage[dvipsnames,svgnames]{xcolor}
\usepackage[shortlabels]{enumitem}
\usepackage[framemethod=TikZ]{mdframed}
\usepackage{amsmath,amssymb,amsthm}
\usepackage[colorlinks]{hyperref}
\usepackage{mathrsfs}
\usepackage{tikz-cd}

\hypersetup{
	colorlinks=true,
	linkcolor=blue,
	filecolor=blue,      
	urlcolor=blue,
	pdftitle={MAT535 Notes}
}

\usepackage{mathtools}
\usepackage{thmtools}
\usepackage[textsize=scriptsize]{todonotes}

\renewcommand{\tilde}{\widetilde}
\renewcommand{\hat}{\widehat}
\renewcommand{\setminus}{\!-\!}
\providecommand{\ol}{\overline}
\providecommand{\eps}{\varepsilon}
\providecommand{\half}{\frac{1}{2}}
\providecommand{\dang}{\measuredangle} %% Directed angle
\providecommand{\dd}{\mathrm d}
\providecommand{\CC}{\mathbb C}
\providecommand{\FF}{\mathbb F}
\providecommand{\HH}{\mathbb H}
\providecommand{\NN}{\mathbb N}
\providecommand{\QQ}{\mathbb Q}
\providecommand{\RR}{\mathbb R}
\providecommand{\ZZ}{\mathbb Z}
\providecommand{\RP}{\mathbb RP}
\providecommand{\CP}{\mathbb CP}
\providecommand{\dg}{^\circ}
\providecommand{\ind}[1]{\mathbf 1_{#1}}
\providecommand{\ii}{\item}
\providecommand{\alert}[1]{{\sffamily\textbf{\textcolor{magenta}{#1}}}}
\providecommand{\opname}{\operatorname}
\providecommand{\li}{\opname{li}}
\providecommand{\ls}{\opname{ls}}
\providecommand{\supp}{\opname{supp}}
\providecommand{\inv}{^{-1}}
\providecommand{\setdiff}{\smallsetminus}
\providecommand{\mb}[1]{\mathbf{#1}}
\providecommand{\abs}[1]{\left\lvert{#1}\right\rvert}
\providecommand{\norm}[1]{\left\|{#1}\right\|}
\providecommand{\dif}{\;d}

\reversemarginpar
\setlength{\marginparwidth}{1in}

\mdfdefinestyle{mdbluebox}{roundcorner=10pt,innerbottommargin=9pt,
	linecolor=blue,backgroundcolor=TealBlue!5,}
\declaretheoremstyle[headfont=\sffamily\bfseries\color{MidnightBlue},
mdframed={style=mdbluebox},]{thmbluebox}

\mdfdefinestyle{mdredbox}{frametitlefont=\bfseries,innerbottommargin=8pt,
	nobreak=true,backgroundcolor=Salmon!5,linecolor=RawSienna,}
\declaretheoremstyle[headfont=\bfseries\color{RawSienna},
mdframed={style=mdredbox},headpunct={\\[3pt]},postheadspace=0pt,]{thmredbox}

\mdfdefinestyle{mdgreenbox}{linecolor=ForestGreen,backgroundcolor=ForestGreen!5,
	linewidth=2pt,rightline=false,leftline=true,topline=false,bottomline=false,}
\declaretheoremstyle[headfont=\bfseries\sffamily\color{ForestGreen!70!black},
mdframed={style=mdgreenbox},headpunct={ --- },]{thmgreenbox}

\mdfdefinestyle{mdorangebox}{linecolor=RedOrange,backgroundcolor=RedOrange!5,
	linewidth=2pt,rightline=false,leftline=true,topline=false,bottomline=false,}
\declaretheoremstyle[headfont=\bfseries\sffamily\color{RedOrange!70!black},
mdframed={style=mdorangebox},headpunct={ --- },]{thmorangebox}

\mdfdefinestyle{mdblackbox}{linecolor=black,backgroundcolor=RedViolet!5!gray!5,
	linewidth=3pt,rightline=false,leftline=true,topline=false,bottomline=false,}
\declaretheoremstyle[mdframed={style=mdblackbox}]{thmblackbox}

\declaretheoremstyle[spaceabove=6pt,spacebelow=6pt]{boldhead}

\declaretheorem[style=thmbluebox,name=Theorem,numberwithin=section]{theorem}

\declaretheorem[style=boldhead,name=Example,sibling=theorem]{example}

\declaretheorem[style=thmbluebox,name=Corollary,sibling=theorem]{corollary}
\declaretheorem[style=thmbluebox,name=Proposition,sibling=theorem]{prop}
\declaretheorem[style=thmbluebox,name=Lemma,sibling=theorem]{lemma}
\declaretheorem[style=thmbluebox,name=Theorem,numbered=no]{theorem*}
\declaretheorem[style=thmbluebox,name=Lemma,numbered=no]{lemma*}
\declaretheorem[style=thmgreenbox,name=Claim,numberwithin=theorem]{claim}
\declaretheorem[style=thmgreenbox,name=Claim,numbered=no]{claim*}
\declaretheorem[style=boldhead,name=Remark,sibling=theorem]{remark}
\declaretheorem[style=boldhead,name=Remark,numbered=no]{remark*}
\declaretheorem[style=thmgreenbox,name=Definition,sibling=theorem]{definition}
\declaretheorem[style=thmgreenbox,name=Definition,numbered=no]{definition*}
\declaretheorem[style=thmorangebox,name=Warning,sibling=theorem]{warning}
\declaretheorem[style=thmorangebox,name=Warning,numbered=no]{warning*}
\declaretheorem[style=thmorangebox,name=Note,numbered=no]{note}
\declaretheorem[style=thmblackbox,name=Exercise,sibling=theorem]{exercise}
\declaretheorem[style=thmblackbox,name=Exercise,numbered=no]{exercise*}

\newenvironment{soln}{\begin{proof}[Solution]}{\end{proof}}
\newenvironment{walkthrough}{\noindent \textbf{Walkthrough.}}{\medskip}
\newlist{walk}{enumerate}{3}
\setlist[walk]{label=\bfseries (\alph*)}

\providecommand{\pow}{\operatorname{Pow}}
%\providecommand{\vocab}[1]{{{\sffamily\textolor{ForestGreen}{#1}}}}
\providecommand{\vocab}[1]{{{\sffamily\textit{#1}}}}
\providecommand{\lcm}{\operatorname{lcm}}
\providecommand{\ord}{\operatorname{ord}}
\providecommand{\Int}{\operatorname{Int}}
\providecommand{\rk}{\operatorname{rk}}
\providecommand{\Mat}{\operatorname{Mat}}
\providecommand{\Ext}{\operatorname{Ext}}
\providecommand{\Tor}{\operatorname{Tor}}

\usepackage{asymptote}
\begin{asydef}
	defaultpen(fontsize(10pt));
	unitsize(1cm);
	size(8cm); // set a reasonable default
	usepackage("amsmath");
	usepackage("amssymb");
	settings.tex="pdflatex";
	settings.outformat="pdf";
	// Replacement for olympiad+cse5 which is not standard
	import geometry;
	// recalibrate fill and filldraw for conics
	void filldraw(picture pic = currentpicture, conic g, pen fillpen=defaultpen, pen drawpen=defaultpen)
	{ filldraw(pic, (path) g, fillpen, drawpen); }
	void fill(picture pic = currentpicture, conic g, pen p=defaultpen)
	{ filldraw(pic, (path) g, p); }
	// some geometry
	pair foot(pair P, pair A, pair B) { return foot(triangle(A,B,P).VC); }
	pair orthocenter(pair A, pair B, pair C) { return orthocentercenter(A,B,C); }
	pair centroid(pair A, pair B, pair C) { return (A+B+C)/3; }
	// cse5 abbreviations
	path CP(pair P, pair A) { return circle(P, abs(A-P)); }
	path CR(pair P, real r) { return circle(P, r); }
	pair IP(path p, path q) { return intersectionpoints(p,q)[0]; }
	pair OP(path p, path q) { return intersectionpoints(p,q)[1]; }
	path Line(pair A, pair B, real a=0.6, real b=a) { return (a*(A-B)+A)--(b*(B-A)+B); }
	// cse5 more useful functions
	picture CC() {
		picture p=rotate(0)*currentpicture;
		currentpicture.erase();
		return p;
	}
	pair MP(Label s, pair A, pair B = plain.S, pen p = defaultpen) {
		Label L = s;
		L.s = "$"+s.s+"$";
		label(L, A, B, p);
		return A;
	}
	pair Drawing(Label s = "", pair A, pair B = plain.S, pen p = defaultpen) {
		dot(MP(s, A, B, p), p);
		return A;
	}
	path Drawing(path g, pen p = defaultpen, arrowbar ar = None) {
		draw(g, p, ar);
		return g;
	}
\end{asydef}

\newcommand{\scr}[1]{\mathscr{#1}}
\newcommand{\obj}{\opname{obj}}
\newcommand{\Ch}{\opname{Ch}}
\renewcommand{\hom}{\opname{Hom}}
\newcommand{\id}{\opname{id}}
\newcommand{\im}{\opname{im}}
\newcommand{\coker}{\opname{coker}}
\newcommand{\tensor}{\otimes}
\newcommand{\dirsum}{\oplus}
\newcommand{\bigtensor}{\bigotimes}
\newcommand{\bigdirsum}{\bigoplus}
\newcommand{\longto}{\longrightarrow}

\title{\vspace{-2.0cm}Ext and Tor}
\author{\large Aditya Pahuja}
\date{\large \today}
\begin{document}
\maketitle
\tableofcontents
\pagebreak

Let $R$ be an (associative, unital) ring.
For a cochain complex $\scr C = (C^n,d^n)_{n\ge 0}$ of $R$-modules,
the idea of the Ext and Tor functors is to relate
the cohomology of $\scr C$ to that of
$\hom_R(\scr C, M)$, $\hom_R(M,\scr C)$, and $\scr C\tensor M$
for some given $R$-module $M$.

\section{Ext}
\begin{definition}[Projective and free resolutions]
    \label{def:resolutions}
    Let $A$ be an $R$-module. A \vocab{projective resolution} of $A$
    is an exact sequence
    \begin{equation*}
	\cdots\longrightarrow P_3\overset{d_3}\longrightarrow P_2\overset{d_2}\longrightarrow P_1
	\overset{d_1}\longrightarrow P_0\overset{d_0}\longrightarrow A\longrightarrow 0.
    \end{equation*}
    in which the $P_n$ are all projective.
    Similarly, a \vocab{free} resolution is such an exact sequence
    where the $P_n$ are all free.
\end{definition}

Every $R$-module has a free (hence projective) resolution:
let $P_0$ be the free module on $A$ (and $d_0$ the obvious projection),
and for $n\ge 1$ let $P_n$ be the free module on $\ker d_{n-1}$
with $d_n$ being the obvious projection onto $\ker d_{n-1}$.
This is often unwieldy; if one has more information about $A$,
there are often nicer resoultions one can construct.
\begin{example}
    If $A$ is projective, then it has a projective resolution
    \begin{equation*}
	0\longto A\longto A\longto 0.
    \end{equation*}
\end{example}

\begin{example}
    Suppose $R$ is a PID and $A = R^n \dirsum\bigdirsum_{i=1}^m R/(p_i^{e_i})$
    is a finitely generated $R$-module.
    There is a projection $p\colon R^{m+n}\to A$
    with kernel $\bigdirsum_{i=1}^m p_i^{e_i} R\cong R^{m}$,
    so that $A$ has a projective resolution
    \begin{equation*}
	0\longto R^{m}\longto R^{m+n}\overset p\longto A\longto 0.
    \end{equation*}
\end{example}

For a given projective resolution $P_n$ of $A$,
apply the contravariant Hom functor $\hom_R(-,B)$
with a fixed $R$-module $B$ to get the cochain complex
\begin{equation*}
    0\longto \hom_R(A,B)\overset{d_0}\longto\hom_R(P_0,B)
    \overset{d_1}\longto\hom(P_1,B)\overset{d_2}\longto\cdots,
\end{equation*}
where the morphisms are the ones that the functor induces
from the morphisms of the resolution.

\begin{definition}[Ext]
    \label{def:ext-def}
    The module $\Ext_R^n(A,B)$ is the cohomology
    of the truncated cochain complex
    \begin{equation*}
	0\longto\hom_R(P_0,B)
	\overset{d_1}\longto\hom(P_1,B)\overset{d_2}\longto\cdots,
    \end{equation*}
    at $\hom_R(P_n,B)$.
\end{definition}

\begin{example}
    Since $\hom_R(-,B)$ is left-exact,
    $\Ext_R^0(A,B) = \ker d_1 = \im d_0 \cong \hom_R(A,B)$.
\end{example}

\begin{example}
    For the free resolution of $\ZZ/n\ZZ$
    \begin{equation*}
	0\longto \ZZ\overset{\cdot n}\longto \ZZ\longto \ZZ/n\ZZ\longto 0
    \end{equation*}
    and an arbitrary abelian group $B$,
    the corresponding cochain complex is
    \begin{equation*}
	0\longto\underbrace{\hom(\ZZ,B)}_{\cong B}
	\overset{\cdot n}\longto\underbrace{\hom(\ZZ,B)}_{\cong B}\longto 0,
    \end{equation*}
    so the Ext groups are
    \begin{equation*}
        \Ext_\ZZ^n(\ZZ/n\ZZ,B) = \begin{cases}
	    {_n}B & n = 0\\
	    B/nB & n = 1\\
	    0 & n \ge 2
        \end{cases}
    \end{equation*}
    where ${_n}B$ is the set of $n$-torsion elements of $B$
    (since each element of the Hom set is uniquely determined
    by the image of $1\in\ZZ/n\ZZ$, which can map to any $n$-torsion element).
\end{example}

Of course, the definition of $\Ext_R^n(A,B)$
a priori requires a choice of resolution,
but it turns out (as the notation suggests)
that it is independent of the resolution.

\todo{pf}

Besides computing projective resolutions,
the other main computational tool for Ext modules
is the long exact sequence in cohomology:
\begin{theorem}[Ext long exact sequence]
    \label{thm:ext-long-exact-sequence}
    Let $0\to L\to M\to N\to 0$ be a short exact sequence of $R$-modules.
    There is an exact sequence
    \begin{align*}
	0\to&\hom_R(N,B)\to\hom_R(M,B)\to\hom_R(L,B)\\
	\overset{\delta_0}\longto
	&\Ext^1_R(N,B)\to\Ext^1_R(M,B)\to\Ext^1_R(L,B)
	\overset{\delta_1}\longto\cdots
    \end{align*}
\end{theorem}

\begin{prop}
    \label{prop:injectivity-criterion}
    For any $R$-module $Q$ the following are equivalent:
    \begin{enumerate}[(1)]
        \ii $Q$ is injective;
	\ii $\Ext_R^1(A,Q) = 0$ for all $R$-modules $A$; and
	\ii $\Ext_R^1(A,Q) = 0$ for all $R$-modules $A$ and all $n\ge 1$.
    \end{enumerate}
\end{prop}

A similar construction can be done with the covariant Hom functor,
but we need a different kind of resolution,
since projective resolutions don't allow use to take advantage
of left-exactness here. Therefore, the obvious choice is to dualize and
take an \vocab{injective resolution};
as the name suggests, this is an exact sequence
\begin{equation*}
    0\longto A\overset{d_0}\longto Q_0\overset{d_1}\longto Q_1
    \overset{d_2}\longto Q_2\overset{d_3}\longto Q_3\longto\cdots
\end{equation*}
in which the $Q_n$ are all injective $R$-modules.









\end{document}
